\documentclass[11pt]{article}
\usepackage[margin=1in]{geometry}
\usepackage{amsmath,amssymb}
\usepackage{enumitem}
\usepackage{bookmark}
\hypersetup{
  pdftitle={PSET 0: Gradescope practice assignment},
  hidelinks}

\title{PSET 0: Gradescope practice assignment}
\author{}
\date{\vspace{-2.5em}}

% Problems are numbered 1-3 across the document.
\newcounter{problem}
\newcommand{\problem}[1]{\stepcounter{problem}\section*{Problem \theproblem: #1}}

\begin{document}
\maketitle

\section*{Purpose}

This ungraded practice assignment is for getting comfortable with
Gradescope before PSET 1 is due. The math is intentionally easy---the
point is the \emph{submission process}, not the content. When you
upload, Gradescope will ask you to \textbf{assign pages to each
question}. Please do this carefully: put each answer on its own page
(or clearly separated), and match each question to the correct page(s)
of your PDF. Submissions can be revised and re-uploaded as many times
as you like before the deadline.

\problem{Simplify expressions}

Simplify the following expressions as much as possible:

\begin{enumerate}[label=\alph*.]
\item \(x^2 x^3\)
\item \(\dfrac{y^5}{y^2}\)
\item \((2a)^3\)
\end{enumerate}

\problem{Solve a linear equation}

Solve for \(x\):
\[
3x + 7 = 22
\]

\problem{Graph sketching}

Sketch the graph of \(y = 2x - 1\), labeling the \(x\)- and
\(y\)-intercepts. (A photo of a hand-drawn sketch is fine---this is
good practice for including images or scanned pages in your Gradescope
submission.)

% Shared AI and Resources statement.
% Include in any problem set with:  % Shared AI and Resources statement.
% Include in any problem set with:  % Shared AI and Resources statement.
% Include in any problem set with:  \input{ai-statement}
\section*{AI and Resources statement}

Please list (in detail) all resources you used for this assignment. If
you worked with people, list them here as well. It is not enough to say
that you used a resource for help, you need to be specific on the link
and \emph{how} it was helpful. W/R/T gen AI tools (including GPT,
Claude, etc.) you cannot use them to do work on your behalf---you
cannot put in any of the questions, etc. You can ask for help on logic
/ sample problems. If you do use GPT or other AI tools, you need to
provide a link to your chat transcript. Any suspected academic
integrity violations will be immediately reported to the Dean of
Students.

\section*{AI and Resources statement}

Please list (in detail) all resources you used for this assignment. If
you worked with people, list them here as well. It is not enough to say
that you used a resource for help, you need to be specific on the link
and \emph{how} it was helpful. W/R/T gen AI tools (including GPT,
Claude, etc.) you cannot use them to do work on your behalf---you
cannot put in any of the questions, etc. You can ask for help on logic
/ sample problems. If you do use GPT or other AI tools, you need to
provide a link to your chat transcript. Any suspected academic
integrity violations will be immediately reported to the Dean of
Students.

\section*{AI and Resources statement}

Please list (in detail) all resources you used for this assignment. If
you worked with people, list them here as well. It is not enough to say
that you used a resource for help, you need to be specific on the link
and \emph{how} it was helpful. W/R/T gen AI tools (including GPT,
Claude, etc.) you cannot use them to do work on your behalf---you
cannot put in any of the questions, etc. You can ask for help on logic
/ sample problems. If you do use GPT or other AI tools, you need to
provide a link to your chat transcript. Any suspected academic
integrity violations will be immediately reported to the Dean of
Students.


\end{document}
